% !TEX root = ../Main.tex
\chapter{课后练习}

\ex[练习 1（多选）]{某次数学考试的一道多项选择题，要求是："在每小题给出的四个选项中，全部选对的得 $5$ 分，部分选对的得 $3$ 分，有选错的得 $0$ 分。"已知该选择题的正确答案是 $CD$，且甲、乙、丙、丁四位同学都不会做（随机猜），下列表述正确的是（\quad）。

\quad A.\ 甲同学仅随机选一个选项，能得 $3$ 分的概率是 $\dfrac{1}{2}$；

\quad B.\ 乙同学仅随机选两个选项，能得 $5$ 分的概率是 $\dfrac{1}{6}$；

\quad C.\ 丙同学随机选择选项，能得分的概率是 $\dfrac{1}{5}$；

\quad D.\ 丁同学随机至少选择两个选项，能得分的概率是 $\dfrac{1}{10}$。}

\sol{
    正确答案 $CD$。"得分"指得 $3$ 分或 $5$ 分，即所选恰是 $\{C\},\{D\},\{C,D\}$ 之一（选了正确答案、且\textbf{没选错}）；只要选了 $A$ 或 $B$ 就得 $0$ 分。逐项判断：

    \textbf{A}：仅选一个，$4$ 种等可能。得 $3$ 分要求"部分选对"，即选中 $C$ 或 $D$，共 $2$ 种，$P=\dfrac{2}{4}=\dfrac{1}{2}$。\textbf{正确}。

    \textbf{B}：仅选两个，从 $4$ 个里选 $2$ 个共 $C_4^2=6$ 种。得 $5$ 分要全选对，即恰好选中 $\{C,D\}$，$1$ 种，$P=\dfrac{1}{6}$。\textbf{正确}。

    \textbf{C}：随机选择（每个选项选或不选，至少选一个），共 $2^4-1=15$ 种。能得分的是 $\{C\},\{D\},\{C,D\}$ 共 $3$ 种，$P=\dfrac{3}{15}=\dfrac{1}{5}$。\textbf{正确}。

    \textbf{D}：至少选两个，选法有 $C_4^2+C_4^3+C_4^4=6+4+1=11$ 种。能得分的只有 $\{C,D\}$ 这 $1$ 种（含 $3$ 个以上必带错选），$P=\dfrac{1}{11}\ne\dfrac{1}{10}$。\textbf{错误}。

    选 \textbf{ABC}。
}

\ex[练习 2]{NBA 总决赛采用 $7$ 场 $4$ 胜制，$2018$ 年总决赛两支球队分别为勇士和骑士。假设每场比赛勇士获胜的概率为 $0.6$、骑士获胜的概率为 $0.4$，且每场结果相互独立，则恰好 $5$ 场比赛决出总冠军的概率为\underline{\hspace{3em}}。}

\sol{
    "恰好 $5$ 场决出冠军"即某队在第 $5$ 场拿到第 $4$ 胜：前 $4$ 场该队 $3$ 胜 $1$ 负，第 $5$ 场该队再胜。两支队分别算、再相加（互斥）。

    勇士夺冠：前 $4$ 场 $3$ 胜 $1$ 负有 $C_4^3$ 种排法，再乘第 $5$ 场胜：
    \[
    C_4^3(0.6)^3(0.4)\cdot 0.6=4\times 0.216\times 0.4\times 0.6=0.20736.
    \]
    骑士夺冠：
    \[
    C_4^3(0.4)^3(0.6)\cdot 0.4=4\times 0.064\times 0.6\times 0.4=0.06144.
    \]
    相加：
    \[
    P=0.20736+0.06144=0.2688=\frac{168}{625}.
    \]
}

\ex[练习 3]{数据 $x_1,x_2,\cdots,x_8$ 的均值为 $\dfrac{5}{2}$，方差为 $2$。现增加一个数据 $x_9$ 后方差不变，则 $x_9$ 的可能取值为\underline{\hspace{3em}}。}

\sol{
    原来 $n=8$，$\bar{x}=\dfrac{5}{2}$，$S^2=2$，由方差简便公式 $S^2=\overline{x^2}-\bar{x}^2$ 反推两个和：
    \[
    \sum_{i=1}^{8}x_i=8\times\frac{5}{2}=20,\qquad
    \sum_{i=1}^{8}x_i^2=8\pr{S^2+\bar{x}^2}=8\pr{2+\frac{25}{4}}=66.
    \]
    加入 $x_9$ 后 $n=9$，新均值 $\bar{x}'=\dfrac{20+x_9}{9}$，新方差仍用简便公式并令其等于 $2$：
    \[
    \frac{66+x_9^2}{9}-\pr{\frac{20+x_9}{9}}^2=2.
    \]
    两边乘 $81$：
    \[
    9(66+x_9^2)-(20+x_9)^2=162.
    \]
    展开整理：
    \[
    594+9x_9^2-400-40x_9-x_9^2=162
    \ \Longrightarrow\
    8x_9^2-40x_9+32=0
    \ \Longrightarrow\
    x_9^2-5x_9+4=0.
    \]
    因式分解 $(x_9-1)(x_9-4)=0$，得
    \[
    x_9=1\quad\text{或}\quad x_9=4.
    \]
}
